% ============================================================================
% COURS 07 — PARTIES 5 À 10
% Reprise du document 11-AcquisitionTraitementSignal2022.docx
% ============================================================================

\section{Représentation spectrale}\label{sec:elec-spectrale}

\subsection{Intérêt de la représentation spectrale}

La représentation temporelle décrit l'évolution d'un signal en fonction du temps. La représentation spectrale décrit la répartition de son contenu en fréquence. Ces deux descriptions sont complémentaires : un même signal peut être observé soit dans le domaine temporel, soit dans le domaine fréquentiel.

La représentation spectrale permet notamment :
\begin{itemize}
  \item d'identifier la valeur moyenne d'un signal ;
  \item de déterminer sa fréquence fondamentale ;
  \item de repérer les harmoniques ;
  \item de connaître la bande fréquentielle occupée ;
  \item de prévoir l'effet d'un filtre sur le signal.
\end{itemize}

\subsection{Dualité temps-fréquence}

À un signal temporel $s(t)$ correspond une représentation fréquentielle $S(f)$. Une variation rapide dans le temps nécessite généralement des composantes de fréquence élevée. À l'inverse, un signal lentement variable possède principalement des composantes de basse fréquence.

Un signal sinusoïdal pur
\[
s(t)=A\cos(2\pi f_0t+\varphi)
\]
possède une seule fréquence $f_0$. Son spectre en amplitude comporte donc une raie à la fréquence $f_0$.

Une composante continue $S_0$ correspond à une raie à la fréquence nulle.

\subsection{Représentation spectrale d'un signal périodique}

Tout signal périodique de période $T$ et de fréquence fondamentale
\[
f_0=\frac{1}{T}
\]
peut, sous certaines conditions, être décomposé en une somme de sinusoïdes : c'est la décomposition en série de Fourier.

On peut écrire :
\[
s(t)=S_0+\sum_{n=1}^{+\infty}A_n\cos\left(2\pi n f_0t+\varphi_n\right).
\]

Dans cette expression :
\begin{itemize}
  \item $S_0$ est la valeur moyenne ;
  \item la composante de fréquence $f_0$ est le fondamental ;
  \item les composantes de fréquence $nf_0$ sont les harmoniques de rang $n$ ;
  \item $A_n$ est l'amplitude de l'harmonique de rang $n$ ;
  \item $\varphi_n$ est sa phase à l'origine.
\end{itemize}

Le spectre d'un signal périodique est un spectre discret : il est constitué de raies séparées de $f_0$.

\subsubsection{Signal sinusoïdal avec valeur moyenne}

Pour
\[
s(t)=S_0+A\cos(2\pi f_0t+\varphi),
\]
le spectre en amplitude contient :
\begin{itemize}
  \item une raie d'amplitude $S_0$ à $f=0$ ;
  \item une raie d'amplitude $A$ à $f=f_0$.
\end{itemize}

\subsubsection{Signal carré symétrique}

Un signal carré symétrique de valeur moyenne nulle contient uniquement des harmoniques impaires. Son développement comporte des termes aux fréquences
\[
f_0,\;3f_0,\;5f_0,\ldots
\]
avec des amplitudes décroissantes lorsque le rang augmente.

\subsubsection{Lecture d'un spectre}

À partir d'un spectre en amplitude, il faut savoir déterminer :
\begin{itemize}
  \item la valeur moyenne ;
  \item la fréquence fondamentale ;
  \item la période du signal ;
  \item le rang de chaque harmonique ;
  \item l'expression temporelle lorsque les phases sont connues ou nulles.
\end{itemize}

\begin{exemple}[Lecture d'un spectre]
Un spectre présente une raie de $\SI{2}{\volt}$ à $f=0$, une raie de $\SI{3}{\volt}$ à $\SI{50}{\hertz}$ et une raie de $\SI{1}{\volt}$ à $\SI{150}{\hertz}$. En supposant les phases nulles :
\[
s(t)=2+3\cos(2\pi 50t)+\cos(2\pi 150t).
\]
La fréquence fondamentale est $\SI{50}{\hertz}$ et la raie à $\SI{150}{\hertz}$ est l'harmonique de rang 3.
\end{exemple}

\subsection{Représentation spectrale d'un signal non périodique}

Un signal non périodique ne possède pas de fréquence fondamentale. Sa représentation fréquentielle est généralement continue : on parle de spectre continu.

Plus un signal est bref dans le temps, plus son spectre est étendu. Cette propriété explique qu'une impulsion très courte nécessite une bande passante importante pour être transmise sans déformation.

\begin{application}
Donner la représentation spectrale en amplitude des signaux temporels proposés dans le document d'origine. Pour chaque signal, préciser la valeur moyenne, le fondamental et le rang des harmoniques visibles.
\end{application}

\begin{application}
À partir des spectres en amplitude fournis, retrouver l'expression temporelle des signaux et indiquer leur valeur moyenne.
\end{application}

% ============================================================================
\section{Filtrage des signaux}\label{sec:elec-filtrage-complet}

\subsection{Définitions}

Un filtre est un système qui modifie le contenu fréquentiel d'un signal. Il laisse passer certaines fréquences et en atténue d'autres.

On distingue :
\begin{itemize}
  \item les filtres passifs, réalisés à partir de résistances, condensateurs et inductances ;
  \item les filtres actifs, qui utilisent en plus des composants actifs, notamment des amplificateurs opérationnels ;
  \item les filtres analogiques, qui agissent sur des signaux continus ;
  \item les filtres numériques, qui agissent sur des suites d'échantillons.
\end{itemize}

\subsection{Comparaison des types de filtre}

Les quatre types fondamentaux sont :
\begin{description}
  \item[Passe-bas :] transmet les basses fréquences et atténue les hautes fréquences.
  \item[Passe-haut :] atténue les basses fréquences et transmet les hautes fréquences.
  \item[Passe-bande :] transmet une bande de fréquences comprise entre deux fréquences de coupure.
  \item[Coupe-bande :] atténue une bande de fréquences et transmet les fréquences situées de part et d'autre.
\end{description}

\subsection{Gabarit d'un filtre}

Le gabarit définit les performances attendues du filtre. Il distingue :
\begin{itemize}
  \item la bande passante, dans laquelle le signal doit être peu atténué ;
  \item la bande atténuée, dans laquelle l'atténuation minimale est imposée ;
  \item la bande de transition, située entre les deux précédentes.
\end{itemize}

Pour un filtre passe-bas, on note généralement :
\begin{itemize}
  \item $f_p$ la limite de la bande passante ;
  \item $f_a$ le début de la bande atténuée ;
  \item $A_p$ l'atténuation maximale admise dans la bande passante ;
  \item $A_a$ l'atténuation minimale exigée dans la bande atténuée.
\end{itemize}

\subsection{Détermination de l'ordre à partir du gabarit}

L'ordre d'un filtre caractérise la rapidité avec laquelle son gain décroît dans la bande de transition. Pour un filtre passe-bas d'ordre $n$, la pente asymptotique est de
\[
-20n\ \text{dB par décade}.
\]

Une estimation de l'ordre minimal peut être obtenue à partir de la pente exigée entre deux points du gabarit :
\[
n\geq
\frac{A_a-A_p}{20\log_{10}(f_a/f_p)}.
\]

L'ordre retenu est l'entier immédiatement supérieur.

\subsection{Décibels}

Le gain en amplitude d'un filtre est
\[
G(f)=\left|\frac{S(f)}{E(f)}\right|.
\]

Le gain en décibels est défini par
\[
G_{\mathrm{dB}}(f)=20\log_{10}G(f).
\]

Quelques valeurs importantes :
\begin{center}
\begin{tabular}{cc}
\toprule
Gain en amplitude & Gain en décibels \\
\midrule
$1$ & $0$ dB \\
$\sqrt{2}$ & $+3$ dB \\
$1/\sqrt{2}$ & $-3$ dB \\
$10$ & $+20$ dB \\
$0{,}1$ & $-20$ dB \\
\bottomrule
\end{tabular}
\end{center}

Pour une puissance, le rapport en décibels s'écrit
\[
G_{\mathrm{dB}}=10\log_{10}\left(\frac{P_s}{P_e}\right).
\]

\subsection{Bande passante et fréquences de coupure}

La fréquence de coupure est généralement définie par une diminution de $\SI{3}{\decibel}$ par rapport au gain maximal. Elle correspond à un gain en amplitude divisé par $\sqrt{2}$.

La bande passante à $-3$ dB est l'ensemble des fréquences pour lesquelles
\[
G_{\mathrm{dB}}\geq G_{\mathrm{dB,max}}-3.
\]

Pour un passe-bande, la largeur de bande vaut
\[
\Delta f=f_{c2}-f_{c1}.
\]

\subsection{Quadripôles}

Un filtre peut être représenté comme un quadripôle possédant :
\begin{itemize}
  \item deux bornes d'entrée ;
  \item deux bornes de sortie ;
  \item une tension d'entrée $u_e(t)$ ;
  \item une tension de sortie $u_s(t)$.
\end{itemize}

Le comportement fréquentiel est décrit par la fonction de transfert complexe
\[
H(j\omega)=\frac{\underline{U_s}}{\underline{U_e}}.
\]

\subsection{Représentation complexe d'un signal sinusoïdal}

À la grandeur sinusoïdale
\[
u(t)=U_m\cos(\omega t+\varphi)
\]
on associe le nombre complexe
\[
\underline{U}=U_m e^{j\varphi}.
\]

La dérivation temporelle correspond à une multiplication par $j\omega$ :
\[
\frac{\mathrm d}{\mathrm dt}\longleftrightarrow j\omega.
\]

L'intégration correspond à une division par $j\omega$.

\subsection{Impédance et admittance complexes}

En régime sinusoïdal, l'impédance complexe d'un dipôle est définie par
\[
\underline{Z}=\frac{\underline{U}}{\underline{I}}.
\]

On écrit aussi
\[
\underline{Z}=R+jX,
\]
où $R$ est la résistance et $X$ la réactance.

L'admittance complexe est
\[
\underline{Y}=\frac{1}{\underline{Z}}=G+jB,
\]
où $G$ est la conductance et $B$ la susceptance.

Les impédances des dipôles élémentaires sont :
\[
\underline{Z_R}=R,
\qquad
\underline{Z_C}=\frac{1}{j\omega C},
\qquad
\underline{Z_L}=j\omega L.
\]

Les lois d'association restent valables avec les impédances complexes :
\begin{itemize}
  \item en série, les impédances s'ajoutent ;
  \item en parallèle, les admittances s'ajoutent.
\end{itemize}

\subsection{Fonction de transfert d'un filtre}

La fonction de transfert est
\[
H(j\omega)=\frac{\underline{U_s}}{\underline{U_e}}.
\]

Son module donne le gain en amplitude :
\[
G(\omega)=|H(j\omega)|.
\]

Son argument donne le déphasage de la sortie par rapport à l'entrée :
\[
\varphi(\omega)=\arg H(j\omega).
\]

\subsection{Méthode de détermination}

Pour déterminer la fonction de transfert d'un filtre passif :
\begin{enumerate}
  \item remplacer chaque dipôle par son impédance complexe ;
  \item repérer les associations série ou parallèle ;
  \item appliquer un diviseur de tension ou les lois de Kirchhoff ;
  \item former le rapport $\underline{U_s}/\underline{U_e}$ ;
  \item simplifier et mettre le résultat sous une forme canonique ;
  \item identifier le gain statique, la pulsation de coupure et l'ordre.
\end{enumerate}

\subsubsection{Filtre RC passe-bas}

\TLFiltreRCPasseBas

\[
H(j\omega)=\frac{1}{1+j\omega RC}
=\frac{1}{1+j\omega/\omega_c},
\qquad
\omega_c=\frac{1}{RC}.
\]

Le module et la phase sont :
\[
|H(j\omega)|=\frac{1}{\sqrt{1+(\omega RC)^2}},
\]
\[
\arg H(j\omega)=-\arctan(\omega RC).
\]

\subsubsection{Filtre RC passe-haut}

\TLFiltreRCPasseHaut

\[
H(j\omega)=\frac{j\omega RC}{1+j\omega RC}.
\]

\subsection{Prévision du comportement aux basses et hautes fréquences}

Une méthode rapide consiste à remplacer les condensateurs et inductances par leurs comportements limites.

Pour un condensateur :
\begin{itemize}
  \item à basse fréquence, $|Z_C|\to +\infty$ : il se comporte comme un circuit ouvert ;
  \item à haute fréquence, $|Z_C|\to 0$ : il se comporte comme un court-circuit.
\end{itemize}

Pour une inductance :
\begin{itemize}
  \item à basse fréquence, $|Z_L|\to 0$ : elle se comporte comme un court-circuit ;
  \item à haute fréquence, $|Z_L|\to +\infty$ : elle se comporte comme un circuit ouvert.
\end{itemize}

Cette étude permet d'identifier rapidement le type de filtre avant tout calcul détaillé.

\subsection{Mise en cascade de cellules élémentaires}

Lorsque plusieurs cellules sont mises en cascade sans interaction entre elles, la fonction de transfert globale est le produit des fonctions de transfert :
\[
H(j\omega)=H_1(j\omega)H_2(j\omega)\cdots H_n(j\omega).
\]

En décibels, les gains s'ajoutent :
\[
G_{\mathrm{dB}}=G_{1,\mathrm{dB}}+G_{2,\mathrm{dB}}+\cdots+G_{n,\mathrm{dB}}.
\]

Les phases s'ajoutent également.

\begin{application}
Donner les fonctions de transfert des filtres proposés dans le document d'origine, puis identifier leur type par l'étude des comportements aux basses et hautes fréquences.
\end{application}

% ============================================================================
\section{Échantillonnage et numérisation}\label{sec:elec-numerisation}

\subsection{Convertisseur analogique-numérique}

Un CAN transforme une tension analogique en un mot binaire. Pour un convertisseur sur $n$ bits, le nombre de codes disponibles est
\[
N=2^n.
\]

Pour une plage de conversion comprise entre $U_{\min}$ et $U_{\max}$, le quantum vaut
\[
q=\frac{U_{\max}-U_{\min}}{2^n}.
\]

Selon la convention retenue dans la documentation, le dénominateur peut être $2^n-1$ lorsque les deux valeurs extrêmes sont représentées exactement. Il faut toujours vérifier la définition donnée par le constructeur.

\subsection{Coder une information sur $n$ bits}

Le code numérique associé à une tension $u_e$ peut être déterminé par
\[
N=\left\lfloor\frac{u_e-U_{\min}}{q}\right\rfloor.
\]

Le résultat doit être limité à l'intervalle des codes admissibles. Une tension hors plage entraîne une saturation.

L'erreur de quantification est liée au remplacement d'une valeur continue par un niveau discret. Elle est généralement comprise entre
\[
-\frac{q}{2}\quad\text{et}\quad +\frac{q}{2}.
\]

\subsection{Technologies des CAN et choix}

Les technologies courantes sont :
\begin{itemize}
  \item CAN flash : très rapide mais nécessitant beaucoup de comparateurs ;
  \item CAN à approximations successives : bon compromis entre vitesse, résolution et coût ;
  \item CAN double rampe : lent mais précis, utilisé dans les appareils de mesure ;
  \item CAN sigma-delta : grande résolution pour des bandes passantes limitées.
\end{itemize}

Le choix dépend notamment :
\begin{itemize}
  \item du nombre de bits ;
  \item de la plage d'entrée ;
  \item du temps de conversion ;
  \item de la fréquence maximale d'échantillonnage ;
  \item de la précision ;
  \item du coût et de la consommation.
\end{itemize}

\subsection{Convertisseur numérique-analogique}

Un CNA transforme un mot binaire en une tension ou un courant analogique. Le quantum et la résolution se déterminent de manière analogue au CAN.

Après conversion, la sortie présente souvent une allure en escalier. Un filtre de reconstruction peut être utilisé pour lisser le signal.

\subsection{Échantillonnage}

L'échantillonnage consiste à prélever la valeur d'un signal analogique à intervalles réguliers $T_e$.

\[
f_e=\frac{1}{T_e}.
\]

\TLEchantillonnageSignal

Une fréquence d'échantillonnage élevée améliore la fidélité de la représentation, mais augmente la quantité de données et impose un temps de conversion plus court.

\subsection{Échantillonneur-bloqueur}

L'échantillonneur-bloqueur prélève la tension d'entrée puis la maintient constante pendant la durée nécessaire à la conversion. Il évite que la tension d'entrée varie pendant le fonctionnement du CAN.

\subsection{Théorème de Shannon-Nyquist}

Pour échantillonner sans ambiguïté un signal dont la fréquence maximale est $f_{\max}$, il faut respecter :
\[
\boxed{f_e>2f_{\max}}.
\]

Dans la pratique, on choisit souvent une fréquence sensiblement supérieure à cette limite.

\subsection{Repliement de spectre}

L'échantillonnage crée des répétitions du spectre autour des multiples de $f_e$. Si $f_e$ est insuffisante, ces répétitions se recouvrent : c'est le repliement de spectre. Des composantes hautes fréquences sont alors interprétées à tort comme des composantes de plus basse fréquence.

\subsection{Filtre anti-repliement}

Le filtre anti-repliement est un filtre passe-bas placé avant le CAN. Il élimine les composantes susceptibles de provoquer un recouvrement spectral.

Son gabarit doit être défini conjointement avec la fréquence d'échantillonnage. À partir du gabarit, on identifie la fréquence maximale utile $f_{\max}$ puis on impose
\[
f_e>2f_{\max}.
\]

\subsection{Déterminer la fréquence d'échantillonnage}

Le choix de $f_e$ doit satisfaire deux contraintes :
\begin{itemize}
  \item la condition de Shannon-Nyquist ;
  \item le temps de conversion du CAN.
\end{itemize}

Si le temps de conversion est $t_c$, il faut disposer d'une période d'échantillonnage compatible :
\[
T_e\geq t_c,
\qquad
f_e\leq \frac{1}{t_c}.
\]

Pour un signal sinusoïdal pur de fréquence $f$, on prend $f_{\max}=f$.

Pour un signal périodique non sinusoïdal, il faut considérer les harmoniques que l'on souhaite conserver. La fréquence maximale utile dépend alors du rang maximal retenu.

% ============================================================================
\section{Codage numérique des informations : numération}\label{sec:elec-numeration}

\subsection{Systèmes de numération}

Dans une base $b$, un nombre s'écrit à l'aide de chiffres compris entre $0$ et $b-1$. Sa valeur décimale est obtenue par
\[
(a_na_{n-1}\ldots a_0)_b=\sum_{k=0}^{n}a_kb^k.
\]

\subsection{Base binaire}

La base 2 utilise les chiffres 0 et 1. Chaque chiffre binaire est appelé bit.

Le bit de poids faible est le LSB et le bit de poids fort le MSB.

Par exemple :
\[
(101101)_2=1\times2^5+0\times2^4+1\times2^3+1\times2^2+0\times2+1=45.
\]

\subsection{Base hexadécimale}

La base 16 utilise les symboles
\[
0,1,\ldots,9,A,B,C,D,E,F.
\]

Chaque chiffre hexadécimal correspond exactement à quatre bits :
\[
(A)_{16}=(1010)_2,
\qquad
(F)_{16}=(1111)_2.
\]

\subsection{Notations}

Les notations courantes sont :
\begin{itemize}
  \item suffixe $_2$ pour le binaire ;
  \item suffixe $_{10}$ pour le décimal ;
  \item suffixe $_{16}$ ou préfixe \,\texttt{0x} pour l'hexadécimal.
\end{itemize}

\subsection{Changement de base}

Pour passer du binaire au décimal, on utilise les puissances de 2.

Pour passer du décimal au binaire, on effectue des divisions successives par 2 et on lit les restes dans l'ordre inverse.

Pour passer du binaire à l'hexadécimal, on regroupe les bits par paquets de quatre à partir de la droite.

\subsection{Addition en base 2}

Les règles élémentaires sont :
\[
0+0=0,
\quad
0+1=1,
\quad
1+1=10,
\quad
1+1+1=11.
\]

Une retenue dépassant le nombre de bits disponibles constitue un débordement.

\subsection{Nombres signés et complément à deux}

Sur $n$ bits, un nombre signé en complément à deux appartient à l'intervalle
\[
-2^{n-1}\leq N\leq 2^{n-1}-1.
\]

Pour coder l'opposé d'un nombre positif :
\begin{enumerate}
  \item écrire sa valeur binaire sur $n$ bits ;
  \item inverser tous les bits ;
  \item ajouter 1.
\end{enumerate}

Le bit de poids fort indique le signe : 0 pour un nombre positif, 1 pour un nombre négatif.

\subsection{Multiplication et division par $2^n$}

Un décalage de $n$ bits vers la gauche correspond, en l'absence de débordement, à une multiplication par $2^n$.

Un décalage de $n$ bits vers la droite correspond à une division entière par $2^n$. Pour les nombres signés, un décalage arithmétique conserve le bit de signe.

\subsection{Autres codes}

D'autres représentations peuvent être utilisées :
\begin{itemize}
  \item code BCD, qui code séparément chaque chiffre décimal ;
  \item code Gray, dans lequel deux valeurs successives ne diffèrent que d'un bit ;
  \item code ASCII pour représenter des caractères.
\end{itemize}

% ============================================================================
\section{Sources}\label{sec:elec-sources}

Cette partie reprend les références et documentations techniques utilisées dans le document d'origine : documentations de capteurs, convertisseurs, filtres et composants électroniques. Les références détaillées seront conservées lors de l'intégration des figures et tableaux d'origine.

% ============================================================================
\section{Exercices du chapitre}\label{sec:elec-exercices}

Les exercices du document source occupent les pages 93 à 135. Ils portent notamment sur :
\begin{itemize}
  \item l'analyse de chaînes d'acquisition ;
  \item le choix et le conditionnement de capteurs ;
  \item l'électrocinétique et les lois de Kirchhoff ;
  \item la représentation spectrale ;
  \item le calcul de fonctions de transfert de filtres ;
  \item les gabarits et l'ordre des filtres ;
  \item les CAN, CNA et l'échantillonnage ;
  \item la numération et le complément à deux.
\end{itemize}

\begin{application}
Les énoncés complets et leurs figures seront intégrés exercice par exercice à partir des pages 93 à 135 du document Word, afin de préserver les schémas et données numériques d'origine.
\end{application}

\section{Synthèse générale}\label{sec:elec-synthese-generale}

\begin{itemize}
  \item Un signal peut être étudié dans les domaines temporel et fréquentiel.
  \item Un signal périodique possède un spectre discret constitué du fondamental et de ses harmoniques.
  \item Un filtre est caractérisé par son gabarit, sa fonction de transfert, son gain et sa phase.
  \item Le comportement limite des condensateurs et inductances permet d'identifier rapidement un filtre.
  \item Le CAN quantifie et code les échantillons d'un signal analogique.
  \item La fréquence d'échantillonnage doit respecter Shannon-Nyquist et le temps de conversion.
  \item Le filtre anti-repliement limite le spectre avant le CAN.
  \item Les bases 2 et 16 et le complément à deux sont indispensables au traitement numérique.
\end{itemize}
